% Support sheet for the meeting (not a report): pictures + short notes.
% Build with:  sh DLR/meshStudy/report/build.sh   (tables and pictures are regenerated from the runs)
% The full report version is in git history (commit 0c8cfe9).
\documentclass[11pt,a4paper]{article}
\usepackage[a4paper,margin=1.8cm]{geometry}
\usepackage[T1]{fontenc}
\usepackage[utf8]{inputenc}
\usepackage{lmodern}
\usepackage{graphicx}
\usepackage{booktabs}
\usepackage{xcolor}
\usepackage{enumitem}
\setlist{itemsep=1pt,topsep=2pt,leftmargin=1.2em}
\graphicspath{{fig/}}
\setlength{\parindent}{0pt}
\setlength{\parskip}{0.5em}
\pagestyle{empty}

\newcommand{\pic}[2]{\IfFileExists{fig/#2}{\includegraphics[width=#1]{#2}}{\fbox{\texttt{\detokenize{#2}} missing}}}
% generated tables carry long report-style captions: drop them here
\newcommand{\nocaptions}{\renewcommand{\caption}[2][]{}}

\begin{document}

{\Large\bfseries Mesh Independence Study, DLR Campus}

\vspace{0.6em}
West wind, 5\,m/s at 10\,m. Steady RANS $k$--$\varepsilon$ (OpenFOAM \texttt{simpleFoam}).
Three meshes: only the cells inside the red box change, \textbf{4\,m $\to$ 2\,m $\to$ 1\,m}.
Everything else (domain, inflow, solver, probes) is identical.

\begin{center}\pic{0.8\textwidth}{box_topview.png}\end{center}
\begin{itemize}
  \item Box around the building with the DLR weather station on it (position known to $\pm$50\,m, the box covers it).
  \item 4 probe lines, each at several heights: \textbf{upstream} (control, should not change),
        \textbf{gap} between buildings, \textbf{wake} behind the building, \textbf{station} above the roof.
\end{itemize}

\newpage
\begin{center}
\pic{0.325\textwidth}{mesh_slice_coarse.png}\hfill
\pic{0.325\textwidth}{mesh_slice_medium.png}\hfill
\pic{0.325\textwidth}{mesh_slice_fine.png}

\vspace{0.4em}
\pic{0.9\textwidth}{mesh_vertical_fine.png}
\end{center}
\begin{itemize}
  \item Top: horizontal cut 7\,m above ground, coarse / medium / fine. Bottom: vertical cut through the station, fine mesh.
  \item 1.19\,M $\to$ 1.28\,M $\to$ 2.07\,M cells. Refining only a box keeps the fine mesh affordable (1\,m everywhere would be tens of millions).
\end{itemize}

\newpage
\begin{center}
\pic{0.325\textwidth}{wind_slice_coarse.png}\hfill
\pic{0.325\textwidth}{wind_slice_medium.png}\hfill
\pic{0.325\textwidth}{wind_slice_fine.png}

\vspace{0.4em}
\pic{\textwidth}{mesh_study_profiles.png}
\end{center}
\begin{itemize}
  \item Same flow pattern on all three meshes; the finer mesh sharpens the slow pockets behind the buildings.
  \item Profiles: \textbf{above $\sim$20\,m the three curves sit on top of each other}; the differences are near the ground between buildings.
  \item Upstream curves are identical: the control behaves.
\end{itemize}

\newpage
\begin{center}
\pic{\textwidth}{mesh_study_convergence.png}

\pic{\textwidth}{residuals.png}
\end{center}
\begin{itemize}
  \item Probe values go flat: every run converged (probe change $<$0.01\,\% over the last 100 iterations).
  \item Fine mesh was run to 2000 iterations because its residuals were still falling at 1000; the probe values did not move.
  \item Medium: residuals plateau higher (p $\approx 8\times10^{-4}$) but its probes are flat.
\end{itemize}

\newpage
{\small\nocaptions
\begin{center}
% generated by make_tex.py
\begin{tabular}{llrrr}
\toprule
Quantity & Usual limit & coarse (4\,m) & medium (2\,m) & fine (1\,m) \\
\midrule
Cell size in the box (m) & -- & 4 & 2 & 1 \\
Total cells & -- & 1\,190\,168 & 1\,275\,741 & 2\,069\,906 \\
Max skewness & $<4$ & \textbf{8.61} & \textbf{8.61} & \textbf{8.61} \\
Faces with skewness $>4$ & 0 & \textbf{5} & \textbf{5} & \textbf{7} \\
Max non-orthogonality (\textdegree) & $<70$ & 64.6 & 64.6 & 64.6 \\
Average non-orthogonality (\textdegree) & -- & 10.6 & 10.3 & 8.4 \\
Faces with non-orthogonality $>70$\textdegree & 0 & 0 & 0 & 0 \\
Max aspect ratio & $<1000$ & 11.5 & 11.5 & 11.5 \\
Min cell volume (m$^3$) & $>0$ & 0.397 & 0.397 & 0.0749 \\
Invalid cells (zero or negative volume) & 0 & 0 & 0 & 0 \\
Wrongly oriented faces (face pyramids) & 0 & 0 & 0 & 0 \\
Failed checkMesh checks & 0 & \textbf{1} & \textbf{1} & \textbf{1} \\
\bottomrule
\end{tabular}


\vspace{0.8em}
% generated by make_tex.py
\begin{tabular}{lrrr}
\toprule
Quantity & coarse (4\,m) & medium (2\,m) & fine (1\,m) \\
\midrule
Status & finished & finished & finished \\
Iterations done & 1000 & 1000 & 2000 \\
Initial residual $U_x$ & $2.4\times10^{-9}$ & $2.4\times10^{-6}$ & $4.9\times10^{-9}$ \\
Initial residual $U_y$ & $2.4\times10^{-8}$ & $5.3\times10^{-5}$ & $2.6\times10^{-8}$ \\
Initial residual $U_z$ & $2.7\times10^{-8}$ & $1.8\times10^{-5}$ & $3.2\times10^{-8}$ \\
Initial residual $p$ & $3.0\times10^{-6}$ & $8.1\times10^{-4}$ & $6.0\times10^{-6}$ \\
Initial residual $k$ & $1.2\times10^{-7}$ & $1.0\times10^{-4}$ & $5.4\times10^{-8}$ \\
Initial residual $\varepsilon$ & $1.1\times10^{-8}$ & $4.8\times10^{-6}$ & $1.1\times10^{-8}$ \\
Continuity error, sum local & $3.1\times10^{-11}$ & $1.8\times10^{-8}$ & $8.2\times10^{-11}$ \\
Continuity error, cumulative & $3.9\times10^{-5}$ & $4.9\times10^{-6}$ & $-1.2\times10^{-5}$ \\
Max probe drift, last 100 it.\ (\%) & $<0.01$ & $<0.01$ & $<0.01$ \\
Wall-clock time (h) & 1.4 & 0.8 & 1.3 \\
\bottomrule
\end{tabular}

\end{center}}
\begin{itemize}
  \item One failed check on every mesh: a few highly skewed faces (5, 5, 7) at building edges, where cells are snapped to the geometry.
        All outside the box on coarse and medium; on fine one is inside the box, 31\,m from the nearest probe.
  \item No invalid cells, non-orthogonality below 70\textdegree.
\end{itemize}

\newpage
{\small\nocaptions% generated by make_tex.py
\begin{table}[htbp]
\centering
\small
\caption{All 19 probes (same columns as Table~\ref{tab:independence}). Heights are above the local ground.}
\label{tab:allprobes}
\begin{tabular}{lrrrrrcr}
\toprule
 & \multicolumn{3}{c}{$|U|$ (m/s)} & \multicolumn{2}{c}{Change (\%)} & & \\
\cmidrule(lr){2-4}\cmidrule(lr){5-6}
Probe (height above ground) & coarse (4\,m) & medium (2\,m) & fine (1\,m) & c.$\to$m. & m.$\to$f. & $p$ & GCI$_\mathrm{fine}$ (\%) \\
\midrule
upstream, 2\,m (control) & 3.12 & -- & -- & -- & -- & -- & -- \\
upstream, 5\,m (control) & 4.34 & -- & -- & -- & -- & -- & -- \\
upstream, 10\,m (control) & 4.98 & -- & -- & -- & -- & -- & -- \\
upstream, 20\,m (control) & 5.94 & -- & -- & -- & -- & -- & -- \\
upstream, 40\,m (control) & 6.55 & -- & -- & -- & -- & -- & -- \\
\midrule
gap, 2\,m & 2.32 & -- & -- & -- & -- & -- & -- \\
gap, 5\,m & 2.39 & -- & -- & -- & -- & -- & -- \\
gap, 10\,m & 2.63 & -- & -- & -- & -- & -- & -- \\
gap, 20\,m & 4.84 & -- & -- & -- & -- & -- & -- \\
gap, 40\,m & 6.61 & -- & -- & -- & -- & -- & -- \\
\midrule
wake, 2\,m & 1.47 & -- & -- & -- & -- & -- & -- \\
wake, 5\,m & 1.91 & -- & -- & -- & -- & -- & -- \\
wake, 10\,m & 2.49 & -- & -- & -- & -- & -- & -- \\
wake, 20\,m & 5.40 & -- & -- & -- & -- & -- & -- \\
wake, 40\,m & 6.54 & -- & -- & -- & -- & -- & -- \\
\midrule
station, 15\,m & 3.72 & -- & -- & -- & -- & -- & -- \\
station, 20\,m & 5.48 & -- & -- & -- & -- & -- & -- \\
station, 25\,m & 5.82 & -- & -- & -- & -- & -- & -- \\
station, 40\,m & 6.59 & -- & -- & -- & -- & -- & -- \\
\bottomrule
\end{tabular}
\end{table}
}
\begin{itemize}
  \item \textbf{20\,m and above: mesh-independent} (2\,m $\to$ 1\,m changes below 1\,\%; gap 20\,m: $-$4.8\,\%).
  \item \textbf{Below $\sim$15\,m between buildings: still changing} by up to 13--15\,\% (about 0.35\,m/s), so 1\,m is not yet proven enough there.
  \item Pressure drop upstream $\to$ wake is tiny (0.17 / 0.46 / 0.72\,Pa against a dynamic pressure of about 15\,Pa), so it is a poor indicator here.
  \item GCI = error band on the fine value; ``osc.'' = values go up then down, so no GCI.
\end{itemize}

\vspace{0.5em}
\textbf{Takeaway.} 2\,m cells are enough for flight above the roofs. Near the ground between buildings,
the result still depends on the mesh at 1\,m.

\textbf{Next.} Repeat medium and fine for the prevailing south-west wind (EDMO: SW 23.5\,\%, W 17.4\,\%);
try a finer box (0.5\,m) or second-order schemes near the ground; then the turbine and the sensor comparison.

\end{document}
